% TODO checklist:
% - [x] Inspected src/mcp/runtime.py (check_permissions function, RBAC integration)
% - [x] Inspected src/security/perm.py (has_perms function, scope checking)
% - [x] Inspected Q&A.md Q17-Q22 for security and RBAC details
% - [x] Inspected src/mcp/tools/security/permissions.py for policy-aware checking
% - [x] Inspected tests/mcp/tools/test_security_permissions.py for authorization tests

\section{Tool Authorization}
\label{sec:tool-authorization}

Tool authorization in the Cineca Agentic Platform implements a defense-in-depth strategy that combines scope-based RBAC, tenant isolation, and policy-aware access control. Every tool invocation passes through multiple authorization checkpoints before execution, ensuring that only authorized principals can perform permitted actions.

\subsection{Authorization Model}
\label{subsec:authorization-model}

The tool authorization model is built on three foundational concepts:

\begin{enumerate}
    \item \textbf{Scopes}: Fine-grained permission tokens (e.g., \texttt{graph:read}, \texttt{admin:all}) that represent specific capabilities.
    \item \textbf{Principals}: Authenticated entities (users, services, machines) with associated scopes and tenant membership.
    \item \textbf{Resources}: The tools and their actions that principals attempt to access.
\end{enumerate}

\subsubsection{Scope Hierarchy}

Scopes follow a hierarchical naming convention that enables both fine-grained and broad permissions:

\begin{lstlisting}[language=Python, caption={Scope hierarchy examples}, label={lst:scope-hierarchy}]
# Fine-grained scopes
"graph:read"        # Read from graph only
"graph:write"       # Write to graph only
"tools:basic"       # Invoke safe tools only

# Broad scopes
"graph:*"           # All graph operations
"tools:all"         # Invoke any tool
"admin:all"         # Full administrative access

# Implied scopes
# "admin:all" implies "tools:all", "graph:*", "security:*", etc.
\end{lstlisting}

Table~\ref{tab:scope-definitions} defines the primary scopes used for tool authorization.

\begin{table}[htbp]
    \centering
    \caption{Primary authorization scopes for MCP tools}
    \label{tab:scope-definitions}
    \begin{tabularx}{\textwidth}{lX}
        \toprule
        \textbf{Scope} & \textbf{Grants Access To} \\
        \midrule
        \texttt{tools:basic} & Safe, read-only tools (catalog, health, status) \\
        \texttt{tools:read} & All read-only tools \\
        \texttt{tools:write} & Tools that modify data \\
        \texttt{tools:all} & All tools except admin \\
        \texttt{tools:admin} & Administrative tools (tenancy, rate limits) \\
        \midrule
        \texttt{graph:read} & Graph query and schema tools \\
        \texttt{graph:write} & Graph CRUD and bulk tools \\
        \texttt{graph:*} & All graph tools \\
        \midrule
        \texttt{security:read} & Permission introspection \\
        \texttt{security:audit} & Audit log access \\
        \texttt{security:admin} & Security configuration \\
        \texttt{security:*} & All security tools \\
        \midrule
        \texttt{model:read} & Model status and testing \\
        \texttt{model:admin} & Model configuration and switching \\
        \midrule
        \texttt{data:read} & Data quality checks \\
        \texttt{data:admin} & Data archival and restoration \\
        \midrule
        \texttt{admin:all} & Full administrative access (implies all scopes) \\
        \bottomrule
    \end{tabularx}
\end{table}

\subsubsection{Principal Extraction}

Principals are extracted from JWT tokens validated by the Auth0 integration. The token payload contains the subject identifier (\texttt{sub}), granted scopes, and tenant membership.

\begin{lstlisting}[language=Python, caption={Principal extraction from JWT payload}, label={lst:principal-extraction}]
# Example JWT payload (decoded)
{
    "sub": "auth0|user_12345",
    "iss": "https://cineca.eu.auth0.com/",
    "aud": ["cineca-agentic-platform"],
    "scope": "openid profile graph:read tools:basic user:me",
    "permissions": ["graph:read", "tools:basic", "user:me"],
    "https://cineca.eu/tenant_id": "tenant_abc",
    "https://cineca.eu/roles": ["researcher"],
    "exp": 1705312800,
    "iat": 1705309200
}
\end{lstlisting}

The MCP runtime normalizes this payload into a principal object:

\begin{lstlisting}[language=Python, caption={Principal normalization in MCP runtime}, label={lst:principal-normalization}]
def _extract_principal_object(principal: Any) -> Any:
    """Normalize principal payloads for permission helpers."""
    if hasattr(principal, "raw"):
        return principal  # Already normalized
    if isinstance(principal, dict):
        return _dict_principal_adapter(principal)
    return principal

class _PrincipalProxy(SimpleNamespace):
    def __init__(self, payload: dict[str, Any]):
        raw_payload = dict(payload)
        sub = raw_payload.get("sub") or raw_payload.get("id")
        scopes = raw_payload.get("scopes") or raw_payload.get("permissions") or []
        super().__init__(
            raw=raw_payload,
            sub=sub,
            scopes=tuple(scopes),
        )
\end{lstlisting}

\subsection{Per-Tool RBAC}
\label{subsec:per-tool-rbac}

Each tool declares its required scope through the \texttt{@mcp\_tool} decorator. The MCP runtime enforces this requirement at invocation time.

\subsubsection{Scope Declaration}

Tools declare their required scope as part of the decorator:

\begin{lstlisting}[language=Python, caption={Tool scope declaration via decorator}, label={lst:tool-scope-declaration}]
@mcp_tool(
    tool_name="graph.query",
    required_scope="graph:read",  # Required scope for this tool
    rate_limit=60,                 # Rate limit per minute
    rate_window=60,                # Rate limit window in seconds
)
async def invoke(ctx: ToolContext, payload: GraphQueryPayload) -> dict:
    """Execute Cypher query against Memgraph."""
    # Implementation...
\end{lstlisting}

\subsubsection{Permission Check Implementation}

The \texttt{check\_permissions} function in \texttt{src/mcp/runtime.py} performs the authorization check:

\begin{lstlisting}[language=Python, caption={Permission checking implementation}, label={lst:check-permissions}]
def check_permissions(
    ctx: ToolContext,
    required_scope: str,
    resource: str | None = None,
    attributes: dict[str, Any] | None = None,
) -> None:
    """
    Check if principal has required permission.
    
    Raises:
        PermissionError_: If permission denied
    """
    from src.security.perm import has_perms
    
    normalized_principal = _extract_principal_object(ctx.principal)
    
    # 1. Verify principal exists
    if not normalized_principal:
        raise PermissionError_(
            "Principal required for this operation",
            details={"required_scope": required_scope}
        )
    
    # 2. Verify tenant match
    if ctx.tenant and hasattr(normalized_principal, "raw"):
        principal_tenant = normalized_principal.raw.get("tenant_id")
        if principal_tenant and principal_tenant != ctx.tenant:
            raise PermissionError_(
                "Tenant mismatch: principal not authorized for this tenant",
                details={"principal_tenant": principal_tenant, "context_tenant": ctx.tenant}
            )
    
    # 3. Check scope using security.permissions module
    if not has_perms(normalized_principal, required_scope):
        raise PermissionError_(
            f"Missing required permission: {required_scope}",
            details={"required_scope": required_scope, "resource": resource}
        )
\end{lstlisting}

\subsubsection{Scope Checking Logic}

The \texttt{has\_perms} function in \texttt{src/security/perm.py} implements the scope matching logic with support for wildcards and implied scopes:

\begin{lstlisting}[language=Python, caption={Scope checking with wildcard support}, label={lst:has-perms}]
def has_perms(principal: Any, required_scope: str) -> bool:
    """
    Check if principal has the required scope.
    
    Supports:
    - Exact match: "graph:read" matches "graph:read"
    - Wildcard: "graph:*" matches "graph:read", "graph:write"
    - Admin override: "admin:all" matches any scope
    """
    if not principal:
        return False
    
    scopes = getattr(principal, "scopes", ()) or ()
    
    # Admin override
    if "admin:all" in scopes:
        return True
    
    # Exact match
    if required_scope in scopes:
        return True
    
    # Wildcard match (e.g., "graph:*" matches "graph:read")
    required_parts = required_scope.split(":")
    for scope in scopes:
        scope_parts = scope.split(":")
        if len(scope_parts) >= 1 and scope_parts[-1] == "*":
            prefix = ":".join(scope_parts[:-1])
            if required_scope.startswith(prefix + ":"):
                return True
    
    return False
\end{lstlisting}

\subsection{Tenant Isolation}
\label{subsec:tenant-isolation}

Multi-tenant isolation is enforced at the tool level to prevent cross-tenant data access.

\subsubsection{Tenant Context Propagation}

The tenant context is propagated through the request lifecycle:

\begin{enumerate}
    \item Client includes \texttt{X-Tenant-ID} header in the request.
    \item The header is validated against the principal's authorized tenants.
    \item The tenant ID is stored in \texttt{ToolContext.tenant}.
    \item Tools use this tenant ID to scope all data operations.
\end{enumerate}

\begin{lstlisting}[language=Python, caption={Tenant validation in authorization flow}, label={lst:tenant-validation}]
# In MCP runtime check_permissions
if ctx.tenant and hasattr(normalized_principal, "raw"):
    principal_tenant = normalized_principal.raw.get("tenant_id")
    
    # Principal must belong to the requested tenant
    if principal_tenant and principal_tenant != ctx.tenant:
        logger.warning(
            "Permission check failed: tenant mismatch",
            extra={
                **ctx.log_context(),
                "principal_tenant": principal_tenant,
                "context_tenant": ctx.tenant,
            }
        )
        raise PermissionError_(
            "Tenant mismatch: principal not authorized for this tenant"
        )
\end{lstlisting}

\subsubsection{Tool-Level Tenant Scoping}

Tools that access data automatically scope queries by tenant:

\begin{lstlisting}[language=Python, caption={Tenant-scoped data access in tools}, label={lst:tenant-scoped-access}]
# In graph.query tool
def _act_run(db: MemgraphAdapter, payload: dict, ctx: ToolContext) -> dict:
    cypher = payload.get("cypher")
    
    # Inject tenant filter for queries that access tenant-scoped data
    if ctx.tenant and not _is_global_query(cypher):
        cypher = _inject_tenant_filter(cypher, ctx.tenant)
    
    results = db.execute(cypher, payload.get("params", {}))
    return {"ok": True, "rows": list(results)}

def _inject_tenant_filter(cypher: str, tenant_id: str) -> str:
    """Add tenant filter to WHERE clause."""
    # Implementation adds WHERE n.tenant_id = $tenant_id
    ...
\end{lstlisting}

\subsection{Policy Files}
\label{subsec:policy-files}

For complex authorization scenarios, the platform supports external policy definitions that can be managed separately from tool code.

\subsubsection{Policy Structure}

Policies are defined in YAML files under \texttt{src/security/policies/}:

\begin{lstlisting}[language=Python, caption={Example policy definition}, label={lst:policy-definition}]
# src/security/policies/tool_access.yaml
policies:
  - name: "graph_researcher_access"
    description: "Researchers can read graph, not write"
    effect: "allow"
    principals:
      roles: ["researcher"]
    actions:
      - "graph:read"
      - "tools:basic"
    resources:
      - "mcp.tools.graph.query"
      - "mcp.tools.graph.schema"
      - "mcp.tools.graph.search"
    conditions:
      tenant_match: true
  
  - name: "admin_full_access"
    description: "Admins have full access"
    effect: "allow"
    principals:
      roles: ["admin"]
    actions:
      - "*"
    resources:
      - "mcp.tools.*"
\end{lstlisting}

\subsubsection{Policy Evaluation}

The \texttt{security.permissions} tool provides runtime policy evaluation:

\begin{lstlisting}[language=Python, caption={Policy evaluation via security.permissions tool}, label={lst:policy-evaluation}]
result = invoke("security.permissions", {
    "action": "evaluate",
    "resource": "mcp.tools.graph.crud",
    "action_type": "write",
    "context": {
        "tenant": "tenant_abc",
        "node_type": "Person"
    }
})
# Response:
{
    "ok": true,
    "allowed": false,
    "policies_applied": ["graph_researcher_access"],
    "denied_by": "graph_researcher_access",
    "reason": "Researchers cannot write to graph"
}
\end{lstlisting}

\subsection{Authorization Decision Flow}
\label{subsec:auth-decision-flow}

Figure~\ref{fig:auth-decision-flow} illustrates the complete authorization decision flow for tool invocations.

\begin{figure}[htbp]
    \centering
    \begin{tikzpicture}[
        node distance=0.6cm and 2cm,
        box/.style={rectangle, draw, rounded corners, minimum width=3cm, minimum height=0.7cm, text centered, font=\small},
        decision/.style={diamond, draw, aspect=2.5, minimum width=2cm, font=\small},
        arrow/.style={->, thick},
        label/.style={font=\scriptsize}
    ]
        % Nodes
        \node[box] (request) {Tool Invocation Request};
        \node[decision, below=of request] (jwt) {JWT Valid?};
        \node[decision, below=of jwt] (principal) {Principal Exists?};
        \node[decision, below=of principal] (tenant) {Tenant Matches?};
        \node[decision, below=of tenant] (scope) {Scope Granted?};
        \node[decision, below=of scope] (policy) {Policy Allows?};
        \node[box, below=of policy] (execute) {Execute Tool};
        
        % Error nodes
        \node[box, right=of jwt] (e401) {401 Unauthorized};
        \node[box, right=of principal] (e403a) {403 No Principal};
        \node[box, right=of tenant] (e403b) {403 Tenant Mismatch};
        \node[box, right=of scope] (e403c) {403 Missing Scope};
        \node[box, right=of policy] (e403d) {403 Policy Denied};
        
        % Arrows
        \draw[arrow] (request) -- (jwt);
        \draw[arrow] (jwt) -- node[left, label] {yes} (principal);
        \draw[arrow] (jwt) -- node[above, label] {no} (e401);
        \draw[arrow] (principal) -- node[left, label] {yes} (tenant);
        \draw[arrow] (principal) -- node[above, label] {no} (e403a);
        \draw[arrow] (tenant) -- node[left, label] {yes} (scope);
        \draw[arrow] (tenant) -- node[above, label] {no} (e403b);
        \draw[arrow] (scope) -- node[left, label] {yes} (policy);
        \draw[arrow] (scope) -- node[above, label] {no} (e403c);
        \draw[arrow] (policy) -- node[left, label] {yes} (execute);
        \draw[arrow] (policy) -- node[above, label] {no} (e403d);
    \end{tikzpicture}
    \caption{Complete authorization decision flow for MCP tool invocations.}
    \label{fig:auth-decision-flow}
\end{figure}

\subsection{Authorization Audit Trail}
\label{subsec:auth-audit}

All authorization decisions are logged for security auditing and debugging.

\paragraph{Successful Authorization.}
\begin{lstlisting}[language=Python, caption={Audit log for successful authorization}, label={lst:auth-success-log}]
logger.debug(
    "Permission check passed",
    extra={
        "tool": "graph.query",
        "action": "run",
        "principal": "auth0|user_12345",
        "tenant": "tenant_abc",
        "trace_id": "trace_xyz",
        "required_scope": "graph:read",
        "granted_scopes": ["graph:read", "tools:basic", "user:me"],
    }
)
\end{lstlisting}

\paragraph{Failed Authorization.}
\begin{lstlisting}[language=Python, caption={Audit log for failed authorization}, label={lst:auth-failure-log}]
logger.warning(
    "Permission check failed: insufficient permissions",
    extra={
        "tool": "graph.crud",
        "action": "create_node",
        "principal": "auth0|user_12345",
        "tenant": "tenant_abc",
        "trace_id": "trace_xyz",
        "required_scope": "graph:write",
        "granted_scopes": ["graph:read", "tools:basic"],
        "error_code": "E_PERMISSION",
    }
)
\end{lstlisting}

\subsection{Testing Authorization}
\label{subsec:testing-auth}

The platform includes comprehensive tests for authorization logic:

\begin{lstlisting}[language=Python, caption={Authorization test examples}, label={lst:auth-tests}]
@pytest.mark.security
def test_tool_requires_scope():
    """Tool invocation without required scope returns 403."""
    # Principal with only basic scope
    principal = {"sub": "user_1", "scopes": ["tools:basic"]}
    ctx = ToolContext(tool="graph.crud", action="create", principal=principal)
    
    with pytest.raises(PermissionError_) as exc_info:
        check_permissions(ctx, required_scope="graph:write")
    
    assert exc_info.value.code == "E_PERMISSION"
    assert "graph:write" in str(exc_info.value)

@pytest.mark.security
def test_admin_scope_grants_all():
    """Admin scope bypasses all permission checks."""
    principal = {"sub": "admin_1", "scopes": ["admin:all"]}
    ctx = ToolContext(tool="tenancy.manage", action="delete", principal=principal)
    
    # Should not raise
    check_permissions(ctx, required_scope="tools:admin")

@pytest.mark.security
def test_tenant_isolation():
    """Cross-tenant access is denied."""
    principal = {"sub": "user_1", "scopes": ["graph:read"], "tenant_id": "tenant_a"}
    ctx = ToolContext(tool="graph.query", action="run", principal=principal, tenant="tenant_b")
    
    with pytest.raises(PermissionError_) as exc_info:
        check_permissions(ctx, required_scope="graph:read")
    
    assert "Tenant mismatch" in str(exc_info.value)
\end{lstlisting}

\subsection{Summary}
\label{subsec:auth-summary}

The tool authorization system provides comprehensive access control through:

\begin{itemize}
    \item \textbf{Scope-based RBAC}: Each tool declares its required scope; the runtime enforces it.
    \item \textbf{Principal validation}: JWT tokens are validated and principals are extracted.
    \item \textbf{Tenant isolation}: Cross-tenant access is prevented at the authorization layer.
    \item \textbf{Policy support}: External policy files enable flexible access rules.
    \item \textbf{Wildcard and implied scopes}: \texttt{admin:all} and \texttt{*} patterns simplify administration.
    \item \textbf{Comprehensive logging}: All authorization decisions are audited.
    \item \textbf{Thorough testing}: Security tests verify authorization behavior.
\end{itemize}

This layered approach ensures that tool access is controlled, auditable, and aligned with enterprise security requirements.
